%%=============================================================================
%% Samenvatting
%%=============================================================================

% TODO: De "abstract" of samenvatting is een kernachtige (~ 1 blz. voor een
% thesis) synthese van het document.
%
% Een goede abstract biedt een kernachtig antwoord op volgende vragen:
%
% 1. Waarover gaat de bachelorproef?
% 2. Waarom heb je er over geschreven?
% 3. Hoe heb je het onderzoek uitgevoerd?
% 4. Wat waren de resultaten? Wat blijkt uit je onderzoek?
% 5. Wat betekenen je resultaten? Wat is de relevantie voor het werkveld?
%
% Daarom bestaat een abstract uit volgende componenten:
%
% - inleiding + kaderen thema
% - probleemstelling
% - (centrale) onderzoeksvraag
% - onderzoeksdoelstelling
% - methodologie
% - resultaten (beperk tot de belangrijkste, relevant voor de onderzoeksvraag)
% - conclusies, aanbevelingen, beperkingen
%
% LET OP! Een samenvatting is GEEN voorwoord!

%%---------- Nederlandse samenvatting -----------------------------------------
%
% TODO: Als je je bachelorproef in het Engels schrijft, moet je eerst een
% Nederlandse samenvatting invoegen. Haal daarvoor onderstaande code uit
% commentaar.
% Wie zijn bachelorproef in het Nederlands schrijft, kan dit negeren, de inhoud
% wordt niet in het document ingevoegd.

\IfLanguageName{english}{%
\selectlanguage{dutch}
\chapter*{Samenvatting}
\lipsum[1-4]
\selectlanguage{english}
}{}

%%---------- Samenvatting -----------------------------------------------------
% De samenvatting in de hoofdtaal van het document

\chapter*{\IfLanguageName{dutch}{Samenvatting}{Abstract}}

Autonome AI-agents verwerken in de Fintech-sector steeds vaker betalingstransacties en andere bedrijfskritische taken. De communicatie met interne systemen verloopt daarbij via het Model Context Protocol (MCP), een standaard die integratie vergemakkelijkt maar tegelijkertijd nieuwe aanvalsvlakken opent: \emph{prompt injection}, \emph{tool poisoning} en \emph{privilege escalation}. Generieke beveiligingstools zoals Burp Suite en Semgrep missen de contextuele kennis om MCP-specifieke kwetsbaarheden systematisch op te sporen, en handmatige code reviews zijn arbeidsintensief en foutgevoelig. Dat 56\% van de bedrijven die AI-agents inzetten beveiligingsincidenten rapporteert in hun MCP-omgevingen, onderstreept de nood aan gespecialiseerde tooling.

De centrale vraag is of een geautomatiseerd security assessment framework MCP-server\-implementaties systematisch kan testen op kritieke kwetsbaarheden, en in hoeverre de effectiviteit verschilt van een gestructureerde handmatige beoordeling. Als evaluatiecase dient een opzettelijk kwetsbare testserver (\emph{Goat}) met vijf bekende, opzettelijk ingebedde kwetsbaarheden.

Het framework combineert statische analyse (SAST) met dynamische fuzzing en draait containergebaseerd via Docker~Compose. Dezelfde Goat-server werd parallel handmatig beoordeeld via code review en HTTP-tests als kwantitatieve benchmark.

Het framework detecteerde alle vijf ingebedde kwetsbaarheden in minder dan 30~seconden, bij een precisie van minimaal 92{,}3\% (12 terechte bevindingen op 13 totaal; één bevinding is debateerbaar als false positive). De handmatige procedure vereiste 47~minuten en identificeerde 4 van de vijf kwetsbaarheden.

De resultaten bevestigen dat geautomatiseerde MCP security assessment haalbaar is en in deze evaluatiecase handmatige beoordeling ruim overtreft in snelheid, bij vergelijkbare of betere recall. De externe validiteit blijft beperkt tot één gecontroleerde case met één reviewer; bredere uitspraken over MCP-servers in het algemeen zijn niet gerechtvaardigd. De voornaamste aanbevelingen zijn: autorisatiechecks verplichten per tool, tool outputs sanitiseren als onvertrouwde invoer, het framework inzetten als screeningslaag en de scanner integreren in CI/CD voor continue compliancemonitoring.
